\section{案例分析}

\subsection{Soang：从搜索到 DPO}

Soang 是一个只能向前推的 box-pushing 游戏，非常考验系统对未来状态的规划能力。Compound AI 的做法拆成两个阶段：先用搜索 trace 学习严格的行动顺序，再用 DPO 让 agent judge 给予反馈并缩小 action spirit。训练流程如下：
\begin{enumerate}
    \item 用 solver 生成完整 trace，记录 cost、plan、intermediate state；
    \item 给 trace 打标签（是否满足 constraints、是否在 4 轮内完成）；
    \item Search Former 学习 trace；agent judge 采样多个 plan、评审 fidelity；
    \item DPO fine-tune 让 agent judge 更偏向最优 plan。
\end{enumerate}
最终结果是：Search Former + DU Forer 用 15M 参数、100K trace 即取得 80\% 成功率，而 solution-only baseline 需要 175M 参数和 1M trace 才能达到 20\%。
\begin{table}[H]
\centering
\begin{tabular}{lrrl}
\toprule
模型 & 参数量 & 样本量 & 成功率 \\
\midrule
Solution-only & 175M & 1M & 20\% \\
Search Former + DU Forer & 15M & 100K & 80\% \\
\bottomrule
\end{tabular}
\caption{Soang 基准的模型对比}
\end{table}

\subsection{Nano optics：从 latent cost 到制造}

在 80\u00d780 网格的光学设计里，搜索 trace 变成了仿真 trace —— agent 观察光线路径、频响和 interference pattern，再把这些信息压入 latent cost \(C\)。整个 pipeline 包括：1）生成 beam splitter、wavelength multiplexer 的仿真；2）挑选满足 brush 约束的 binary pattern；3）用 solver 输出\(X^\star\)；4）将 loss 反向传至 \(C\)。
\begin{importantbox}{从 trace 到制造的闭环}
1）Trace 保留了光线干涉与制造笔刷的信息；2）latent \(C\) 既提取环境描述又控制 solver；3）solver 反馈被用来校准 \(C\)，避免生成不可制造的点；4）multi-task loss 让设计在不同 benchmark 之间迁移。
\end{importantbox}
\begin{knowledgebox}{案例启发}
这两个案例共同说明：Compound AI 可以把规划 trace、交互问答、Latent solver 串联起来，每个输出既是下一阶段的输入，又是可审计的中间结果。这样的 modular pipeline 更容易定位失败与进行小步迭代。
\end{knowledgebox}

\subsection{本章小结}

案例分析强调了工程实践中的数据流：search trace、agent judge、latent cost 彼此连接，并在不同任务上反复验证，确保抽象理论可以落地。

